% Source: https://artofproblemsolving.com/community/c3h1656531_2018_ms_mte_sprint_all
% Stable identity: aops-topic-1656531-q20
Let $S_n=x_1+\cdots+x_n$. The digits of $1,\dots,9$ sum to $45$. In $10,\dots,69$ each tens digit $1,\dots,6$ appears $10$ times and the units digits $0,\dots,9$ appear $6$ times each, giving $10\cdot21+6\cdot45=480$. So after writing $69$ the sum is $525$. Continuing digit by digit:
\[
70:\ 532,532;\quad 71:\ 539,540;\quad\ldots;\quad 77:\ 602,609;\quad 78:\ 616,624;\quad 79:\ 631,640 .
\]
The sum jumps from $631$ to $640$, so $639$ is never a partial sum.

The other values do occur: the digits of $1,\dots,99$ sum to $900$ and those of $100,\dots,109$ to $55$, so $S$ reaches $955$ after $109$ and $956$ after the first digit of $110$; similarly $1198$ is the sum after writing $142$, $8931$ after writing $740$, and $14503$ after the first digit of $1101$ (found by continuing the count). The answer is $\boxed{\textbf{(a)}\ 639}$.
